% APPROACH HINT: If three points p_1,p_2,p_3 in S∩(S+v) then S contains 6 points p_1,p_2,p_3,p_1-v,p_2-v,p_3-v; use the cyclic ordering of the convex hull to show this is impossible. Concretely: label the vertices of conv(S) cyclically as s_0,s_1,...,s_{n-1}. A point p_i in S∩(S+v) gives a directed chord from p_i-v to p_i with direction vector v. Three such parallel chords in a convex polygon lead to a contradiction via the cyclic order. For the equality case (|S∩(S+v)|=2), show p,q,p-v,q-v are distinct and form a parallelogram by verifying p+(q-v)=q+(p-v). Output must be a self-contained LaTeX fragment with no \documentclass header, just \begin{theorem}...\end{theorem} and \begin{proof}...\end{proof} environments. Include explicit \begin{proof}...\end{proof} blocks for both claims.
\documentclass{article}
\usepackage{amsmath,amssymb,amsthm}
\begin{document}

\textbf{Subproblem 2 --- Two distinct translates meet in at most two points.}

\medskip
\textbf{Statement.}
Let $S\subset\mathbb{R}^2$ be a finite set of points in convex position, and let $v\in\mathbb{R}^2$, $v\neq 0$.
Prove:
\begin{enumerate}
  \item $|S\cap(S+v)| \le 2$.
  \item If $|S\cap(S+v)|=2$, say $S\cap(S+v)=\{p,q\}$, then the four points $p,\, q,\, p-v,\, q-v$ are distinct vertices of $S$ and form a parallelogram (i.e.\ $p-v$ and $q-v$ are also in $S$, and these four points satisfy $(p)+(q-v)=(q)+(p-v)$).
\end{enumerate}
(Equivalently: for distinct $t,t'\in\mathbb{R}^2$, two translated copies $S+t$ and $S+t'$ share at most two points.)

\medskip
\textbf{Hint.}
Order the points of $S$ cyclically along the boundary of $\operatorname{conv}(S)$.  A point $p\in S\cap(S+v)$ means both $p\in S$ and $p-v\in S$, so $v = p-(p-v)$ is the difference of two distinct points of $S$.  If three points $p_1,p_2,p_3\in S\cap(S+v)$ all satisfy $p_i - v\in S$, then $S$ contains three distinct pairs $(p_1, p_1-v)$, $(p_2,p_2-v)$, $(p_3,p_3-v)$ with the same difference $v$.  Use convex position to derive a contradiction.

\end{document}
